\documentclass[a4paper,11pt]{article}

\usepackage[T1]{fontenc}
\usepackage[utf8]{inputenc}
\usepackage{tgtermes}
\usepackage{mathptmx}
\usepackage[english]{babel}
\usepackage[a4paper,top=2.5cm,bottom=2.5cm,left=2.5cm,right=2.5cm]{geometry}
\usepackage{amsmath}
\usepackage{amssymb}
\usepackage{graphicx}
\usepackage{booktabs}
\usepackage[colorinlistoftodos]{todonotes}
\usepackage[colorlinks=true, allcolors=blue]{hyperref}
\usepackage[numbers]{natbib}
\usepackage{setspace}
\doublespacing

\setlength{\marginparwidth}{2cm}
\setcounter{secnumdepth}{2}
\setlength{\parindent}{1.5em}
\setlength{\parskip}{0pt}

\title{\textbf{Bank of China's ``Crude Oil Treasure'' Incident: \\
A Risk Management Case Study}}
\author{Ng Tze Kean}
\date{Mar 2026}

\begin{document}

\maketitle

\section{Introduction}

In April 2020, international crude oil markets experienced an
unprecedented shock when the West Texas Intermediate (WTI) futures
contract for May delivery settled at a deeply negative price of
approximately $-$37.63 USD per barrel. This event, the first time in
history that U.S. crude oil futures had traded below zero, exposed
vulnerabilities across global commodity markets and brought to light
a particularly striking case of institutional risk management failure
at the Bank of China (BoC). The bank's retail investment product
known as ``Crude Oil Treasure'' (Yuan You Bao), which was linked to
WTI and Brent crude futures, became the focal point of an incident
that resulted in estimated losses of 7 to 10 billion yuan, affecting
approximately 60,000 retail investors and triggering significant
regulatory sanctions against the bank \cite{caixin2020stops}.

This report examines the Bank of China Crude Oil Treasure incident
through the lens of bank risk management. The analysis proceeds in
several stages. First, it provides the background context of the
April 2020 oil price crash, explaining the confluence of factors that
drove WTI prices into negative territory. Second, it describes the
product structure and mechanics of Crude Oil Treasure and explains
how both retail clients and the bank itself suffered financial
losses. Third, it analyses the specific risk exposures involved,
including market risk, liquidity risk, credit risk, and operational
risk, with reference to BoC's balance sheet position. Fourth, it
discusses how the bank might have avoided or mitigated the failure
with the benefit of hindsight, drawing on comparisons with peer
institutions and post-crisis regulatory guidance. The report
concludes with reflections on the broader lessons for commercial bank
risk management.

\section{Background: The April 2020 Oil Price Crash}

The collapse of WTI crude oil futures to negative prices on 20 April
2020 was the product of an extraordinary convergence of
macroeconomic, geopolitical, and market-structural forces. Beginning
in early 2020, the rapid global spread of the Covid-19 pandemic
caused an unprecedented contraction in oil demand, particularly from
the transportation and industrial sectors that together account for
the majority of global crude consumption. Simultaneously, a price war
between Saudi Arabia and Russia in March 2020 led major producers to
maintain or increase output at a time when demand was falling
sharply, resulting in a massive and growing surplus of crude oil
worldwide \cite{atlantis2022casestudy}.

The physical infrastructure of the oil market proved unable to absorb
this surplus. Storage capacity at Cushing, Oklahoma---the delivery
point for WTI futures contracts---filled rapidly, reaching more than
70 to 75 percent of working capacity by mid-April, with much of the
remaining capacity already committed under long-term contracts. This
was a critical development because WTI futures are physically settled
instruments: a long holder at expiry must accept delivery of 1,000
barrels per contract, and if storage is unavailable, the holder has
no way to fulfil this obligation without finding a counterparty
willing to take on the physical commodity \cite{wharton2024benchmarks}.

The Chicago Mercantile Exchange (CME) had recognized the emerging
risk. On 15 April 2020, the CME explicitly warned market participants
that energy futures, including WTI, could trade at or below zero due
to the scarcity and cost of storage, and it updated its systems to
accommodate negative pricing. However, many market participants and
intermediaries did not fully internalise this warning or update their
own systems accordingly \cite{cgtn2020fined}.

On 20 April 2020, with the May WTI contract entering its final
trading day, the market dynamics became acutely one-sided. Financial
and retail long holders who could not handle physical delivery
scrambled to exit their positions, but the pool of willing buyers was
severely constrained because storage and logistical capacity at
Cushing was nearly exhausted. As the trading session progressed,
selling pressure intensified while depth in the order book thinned
dramatically. At approximately 2:08 p.m. New York time, WTI futures
broke through zero for the first time in history, and the contract
eventually settled near $-$37.63 USD per barrel. In effect, sellers
were paying buyers to assume the contractual obligation to accept oil
that the sellers themselves could not store \cite{atlantis2021fmet}.

Academic and policy analyses have concluded that the negative
settlement price was primarily driven by binding storage constraints
at Cushing, amplified by a large open interest held by financial
traders and funds that were structurally unable to accept physical
delivery and thus became forced sellers into a severely illiquid
market \cite{wharton2024benchmarks}.

\section{The Crude Oil Treasure Product}

\subsection{Product Structure and Design}

``Crude Oil Treasure'' (Yuan You Bao) was launched by Bank of China
in January 2018 as a wealth-management product designed for mainland
Chinese retail investors, providing them with synthetic exposure to
international crude oil markets through domestic and foreign futures
contracts, primarily WTI and Brent. The product was structured so
that customers paid 100 percent cash margin to BoC, which then
established offsetting positions in overseas crude oil futures on its
own account. Investors were credited or debited daily profit and loss
based on changes in futures prices, without any obligation to take or
make physical delivery of oil \cite{drpress2023icbc}.

The product was registered and marketed as a medium-risk (R3)
wealth-management product. Marketing materials emphasised that the
product required 100 percent margin and involved ``no leverage,''
creating the impression among many investors that their maximum
possible loss was limited to their initial principal investment.
However, the contractual terms and the underlying futures replication
meant that losses could theoretically exceed deposits if the
referenced futures contract settled at deeply negative prices---a
scenario that, prior to April 2020, had never occurred and was not
contemplated in the product's risk disclosures \cite{drpress2023icbc}.

\subsection{How Clients Lost Money}

Bank of China had scheduled the rollover and settlement of its
WTI-linked Crude Oil Treasure positions to take place on the last
trading day of the May contract, which fell on 20 April 2020 U.S.
time. Unlike several peer Chinese banks---notably Industrial and
Commercial Bank of China (ICBC) and China Construction Bank (CCB),
which had rolled their similar paper-crude products to the June
contract several days before expiry---BoC held its positions until
the final trading session, directly exposing its clients to the
historic negative settlement price \cite{atlantis2022casestudy}.

The consequences for investors were devastating. Approximately 60,000
retail clients had collectively deposited around 4.2 billion yuan in
margin. When their positions were settled at the negative WTI price,
not only did they lose the entirety of their margin, but they also
collectively owed an additional 5.8 billion yuan to the bank,
implying total client losses of approximately 10 billion yuan. Many
investors found themselves with deeply negative account balances,
meaning that the bank considered them to owe money beyond what they
had originally invested \cite{argus2020fined}.

Compounding the problem, BoC's trading rules stipulated that retail
positions would be auto-liquidated if margin fell below 20 percent of
the notional exposure. However, the bank's risk and trading systems
had not been designed to handle negative futures prices. As a result,
the auto-liquidation mechanism failed to trigger as the price fell
through zero, allowing client losses to accumulate far beyond the
margin threshold and turning many accounts deeply negative. This
system failure was a direct contributor to the severity of the
losses, as an operational auto-close at zero or any small positive
price would have significantly limited client exposure
\cite{atlantis2022casestudy}.

\subsection{How the Bank Lost Money}

Bank of China's own financial exposure to the incident was
significant, though distinct from the client losses. Economically,
Crude Oil Treasure was structured as an off-balance-sheet derivatives
business in which BoC transformed retail deposits into margin
collateral for its own futures positions, passing market profit and
loss through to customers while retaining execution and counterparty
risk on its own books \cite{bloomberg2020losses}.

Once the settlement produced deeply negative client account balances,
these balances became unsecured credit exposures from the bank to
tens of thousands of retail customers. If those customers refused or
were unable to pay---and many did refuse, on the grounds that the
product had been mis-sold and the auto-liquidation system had
failed---the bank would be forced to absorb the shortfall as a
trading loss that would crystallise on its balance sheet as a credit
loss, accompanied by severe reputational damage \cite{bloomberg2020losses}.

Under intense regulatory and political pressure, BoC ultimately
offered most small investors compensation of up to 20 percent of
their original investment, capped at 10 million yuan per client. This
compensation plan was estimated to cost the bank approximately 6 to
6.5 billion yuan, on top of its own hedging and operational losses
\cite{caixin2020mulls}. The China Banking and Insurance Regulatory
Commission (CBIRC) additionally fined the bank 50.5 million yuan,
sanctioned responsible managers, and suspended certain related
business lines \cite{reuters2020fined}.

\section{Risk Exposure Analysis}

\subsection{Market and Price-Gap Risk}

The primary risk that materialised in the Crude Oil Treasure incident
was extreme market risk. The collapse of WTI from positive territory
to $-$37.63 USD within a single trading session was far beyond any
historical precedent, yet BoC's stress tests and risk limits appear
not to have contemplated negative-price scenarios. Academic analyses
of the incident argue that the product's design failed to fully
reflect the characteristics of futures contracts, including
convergence to spot price near expiry and the risk of severe price
dislocations when liquidity dries up in physically-settled markets
\cite{oilprice2020fine}.

From a commercial-banking perspective, this represents a classic
trading-book market-risk failure. The bank's risk limits and stress
scenarios were anchored on historical volatility rather than
tail-risk events, and its decision to roll and settle positions on
the most illiquid day of the contract cycle maximised exposure to
intraday price gaps. Because BoC's hedge positions in futures
mirrored client exposures, any unhedged gap or execution slippage
around the rollover directly translated into profit-and-loss impacts
for the bank and ultimately consumed its capital \cite{atlantis2021analysis}.

\subsection{Liquidity Risk}

The incident exposed two distinct dimensions of liquidity risk. The
first was market liquidity risk: near expiry, the May WTI order book
thinned dramatically as storage constraints at Cushing limited the
number of willing buyers, leading to violent price swings and periods
where large positions simply could not be unwound at reasonable
prices. Research on the Crude Oil Treasure case highlights that
time-zone differences between Chinese retail trading hours and the
CME market, combined with BoC's suspension of new WTI-linked trading
on 21 April, further impaired the bank's flexibility in managing its
positions and executing orderly liquidation \cite{atlantis2021fmet}.

The second dimension was funding liquidity risk. As WTI prices moved
sharply against BoC's hedging futures positions, the bank would have
faced substantial intraday margin calls from its clearing brokers.
These calls had to be met from the bank's own liquidity buffers or
wholesale funding channels, creating a sudden and large cash-flow
demand on the bank's treasury. While BoC's overall liquidity position
remained adequate given its scale, the episode demonstrated how a
retail derivatives product can unexpectedly generate significant
funding stress, competing with core lending and payment obligations
on the liability side of the balance sheet \cite{joebm2024analysis}.

To place these liquidity demands in the context of BoC's balance
sheet, it is instructive to consider the bank's financial position at
year-end 2020. Bank of China's total assets stood at approximately
24.40 trillion yuan, with equity attributable to shareholders of
roughly 2.04 trillion yuan, implying a balance-sheet leverage ratio
of approximately 12 times---typical for a large commercial bank. The
bank's capital adequacy ratio remained above regulatory minimums.
Even assuming total net losses from the Crude Oil Treasure incident
of 7 to 9 billion yuan including compensation costs, this represents
less than 0.5 percent of shareholder equity, which did not threaten
solvency \cite{atlantis2021analysis}. However, it did consume
economic capital and raised serious questions about the risk-return
trade-off of engaging in non-core trading activities on an already
leveraged commercial-banking balance sheet. The incident illustrates
how even small mis-priced derivative businesses, when distributed
across tens of thousands of retail clients, can generate material
profit-and-loss volatility and capital consumption that is
disproportionate to the fee income they generate.

\subsection{Credit and Counterparty Risk}

The settlement at negative prices transformed what had been marketed
as a fully-margined investment product into an unsecured retail
credit exposure for the bank. According to reporting by Caixin,
clients had deposited about 4.2 billion yuan of margin but, after
settlement, collectively owed an additional 5.8 billion yuan, an
amount that many were unwilling or unable to pay. BoC initially
asserted its contractual rights to claim these negative balances, but
faced both legal challenges from investors and political pressure
from regulators, so much of this exposure effectively became
non-performing or disputed retail credit claims \cite{caixin2020exclusive}.

From a credit-risk management perspective, BoC had effectively
extended unsecured credit lines to tens of thousands of retail
clients without recognising or treating them as such within its
credit portfolio. Because the product was marketed as fully margined
and was not framed internally as a lending exposure, it was not
subject to the credit underwriting standards, concentration limits,
or provisioning requirements that would ordinarily apply to retail
lending. Regulators subsequently instructed banks to cease selling
wealth-management products that could expose investors to unlimited
losses, recognising that such structures blur the boundary between
investment and leveraged speculation and can create unrecognised
retail credit risk for the sponsoring bank \cite{reuters2020fined}.

\subsection{Operational, Model, and Conduct Risk}

The CBIRC investigation concluded that BoC had exhibited ``irregular
product management,'' including unclear contract terms regarding
margin mechanisms and a failure to rectify known defects in trading
system functions, particularly the inability of the auto-liquidation
system to handle negative prices. These deficiencies constitute
operational and model risk: critical systems embedded faulty
assumptions---specifically, that prices could not be negative---and
governance structures failed to ensure that information technology,
legal, and risk functions updated their models and documentation
after the CME signalled the possibility of negative pricing on 15
April 2020 \cite{cgtn2020fined}.

The regulator also found that BoC's internal controls on
consumer-rights protection and sales management were inadequate.
Promotional materials were in some cases exaggerated and did not
clearly explain that investors could lose more than their principal,
despite the R3 ``moderate risk'' classification. This represents a
conduct-risk failure with direct commercial-banking implications:
mis-selling to retail clients damages franchise value and can force
the bank, under regulatory or political pressure, to absorb losses
that are legally the customer's responsibility, thereby converting
operational and reputational risk into real capital costs
\cite{caixin2020probe}.

\section{Hindsight: How Bank of China Could Have Avoided the Failure}

With the benefit of hindsight, several concrete measures could have
prevented or substantially mitigated the Crude Oil Treasure disaster.
These span product design, operational systems, rollover policy, and
risk governance.

\subsection{Adopting an Earlier Rollover Schedule}

Perhaps the single most consequential decision was BoC's choice to
schedule its contract rollover on the last trading day of the May WTI
futures. By contrast, ICBC and CCB rolled their equivalent
paper-crude products to the June contract several days before the May
expiry, thereby avoiding settlement at the extreme negative price
entirely. Had BoC adopted a similar conservative rollover
policy---settling positions at least five to seven business days
before the expiry date---its clients' positions would have been
closed at prices that, while low, were still well within positive
territory. Academic analysis of the incident has identified the late
rollover timing as the single largest controllable factor in the
scale of losses, noting that early rollover is standard practice
among institutional market participants precisely because liquidity
deteriorates sharply in the final days before physical-delivery
expiry \cite{drpress2023icbc}. Regulatory commentary from the CBIRC
also pointed to the rollover timing as a key deficiency in BoC's risk
management framework \cite{caixin2020probe}.

\subsection{Updating Systems for Negative Prices}

The CME issued an advisory on 15 April 2020---five days before the
crash---explicitly warning that energy futures could trade at or
below zero and that market participants should prepare their systems
accordingly. BoC's trading and risk management systems were not
updated to handle negative prices, which caused the auto-liquidation
mechanism to fail. Had the bank responded promptly to the CME
advisory by modifying its systems to accommodate negative pricing and
by activating the 20 percent margin auto-close provision at or near
zero, client losses would have been limited to their margin deposits
rather than extending into deeply negative territory
\cite{cgtn2020fined}. This failure illustrates the importance of
maintaining adaptive operational infrastructure and of treating
exchange-level advisories as urgent triggers for system review,
particularly when they signal the possibility of fundamentally novel
market conditions.

\subsection{Strengthening Product Design and Risk Classification}

The classification of Crude Oil Treasure as an R3 medium-risk product
was fundamentally at odds with its actual payoff profile, which
replicated a futures contract with theoretically unlimited downside.
A more rigorous new-product approval process, with independent review
by risk, legal, and compliance functions, should have identified this
mismatch before the product was launched. Specifically, the product
should either have been reclassified to a higher risk tier with
correspondingly stricter investor suitability requirements, or it
should have incorporated explicit loss-limiting features such as a
contractual floor at zero (so that client losses could never exceed
the margin deposited) or mandatory stop-loss triggers at specified
threshold levels \cite{scmp2020billion}. As the Harvard Business
School case study on the incident notes, the structural mismatch
between what the product's risk label communicated and what its
contractual terms actually permitted was a root cause of both the
mis-selling and the scale of the eventual losses \cite{hbs2020structured}.

\subsection{Implementing Robust Stress Testing}

BoC's stress-testing framework for the Crude Oil Treasure product
evidently did not contemplate the possibility of negative oil prices.
Even before the April 2020 event, the academic literature on
commodity markets had noted that physical-delivery futures could
theoretically settle at negative values when storage is scarce and
costly. A more robust stress-testing regime---one that incorporated
extreme but plausible scenarios including negative settlement prices,
sudden storage capacity exhaustion, and complete evaporation of
market liquidity near expiry---would have revealed the product's
vulnerability to tail-risk events. Such analysis might have prompted
the bank to hedge against extreme downside scenarios, to establish
stricter position limits, or to revise the product's contractual
terms \cite{atlantis2021fmet}. Peer institutions' experience
demonstrates that even simple heuristic risk controls, such as
mandatory early rollover rules, can serve as effective
circuit-breakers against tail-risk materialisation.

\subsection{Improving Investor Suitability and Disclosure}

Finally, the conduct and suitability failings identified by the CBIRC
could have been addressed through clearer and more conservative
investor disclosures. If marketing materials had explicitly stated
that losses could exceed 100 percent of principal under extreme
market conditions, and if the investor qualification process had
restricted access to the product to clients with demonstrated
understanding of derivatives risk, both the number of affected
investors and the intensity of the subsequent regulatory and
reputational fallout would likely have been reduced
\cite{caixin2020bitter}. Post-crisis regulatory guidance has since
mandated these kinds of enhanced disclosures for derivative-linked
wealth-management products, reflecting the lessons of the Crude Oil
Treasure episode \cite{reuters2020fined}.

\section{Conclusion}

The Bank of China Crude Oil Treasure incident stands as a cautionary
example of how failures in product design, risk management,
operational systems, and investor protection can converge to produce
outsized losses in a commercial-banking setting. The April 2020
collapse of WTI oil prices to negative territory was an extreme and
unprecedented market event, but its impact on BoC and its clients was
magnified by a series of avoidable decisions: a late rollover
schedule that maximised exposure to illiquid end-of-contract trading,
systems that could not accommodate negative prices, a risk
classification that understated the product's true downside, and
disclosures that left retail investors unaware of their potential for
losses beyond their principal.

The incident consumed an estimated 7 to 9 billion yuan of economic
value, imposed regulatory sanctions and reputational damage on the
bank, and prompted a comprehensive tightening of regulatory standards
for derivative-linked wealth-management products in China. While the
losses represented less than 0.5 percent of Bank of China's
shareholder equity and did not threaten the bank's solvency, they
demonstrated that on a leveraged commercial-banking balance sheet,
even seemingly small non-core trading activities can generate
material capital consumption and profit-and-loss volatility that is
disproportionate to the fee income they produce.

For bank managers, the central lesson is that any product whose
economics replicate those of a futures or options contract must be
governed with the same rigour as a derivatives business, regardless
of how it is labelled or marketed. This requires robust stress
testing that contemplates extreme scenarios, adaptive operational
systems that can respond to changing market conditions, conservative
rollover and liquidity management policies, honest risk
classification, and clear investor disclosures. The Crude Oil
Treasure incident shows that when these safeguards are absent, the
consequences can be severe---not only for the bank's clients, but for
the institution's balance sheet, reputation, and regulatory standing.

\newpage

\bibliographystyle{plainnat}
\bibliography{references}

\end{document}
